\scanpage{403}\mspage{142}
Le cas qui nous occupe ici serait celui où
$\mathbb D=\mathrm{HOT}$, et on prend $H^{\mathbb D}_{\bullet}$, et $a={}$objet final de
$\mathbb D(e)=\mathrm{Hot}$.)

$5^\circ$) Si $1^\circ,2^\circ,3^\circ$ sont vérifiés, on est moralement
sûr que la construction de $\underline{\mathrm{Ch}}_{\infty}(X)$ est bel
et bien l'équivalent exact du formalisme de
Cartan-Serre. Pour en avoir le cœur totalement
\unsure{net}, i.e. raccrocher \inserted{explicitement} cette théorie aux
constructions classiques, il faudrait \deleted{\unsure{poser}} prouver un
\uline{énoncé de comparaison} avec la théorie
semisimpliciale, qui dirait essentiellement
ceci : Si on désigne, pour $X$ dans $\mathrm{Cat}$, par
$\widetilde X$ l'objet associé de $\mathbb{\Delta}^{\wedge}$ (le ``nerf'' de $X$),
alors la flèche
\diagram{\begin{tikzcd}[ampersand replacement=\&,row sep=tiny]
\widetilde{\underline{\mathrm{Ch}}_{\infty}(X)} \arrow[r] \& \widetilde{X\times X} \arrow[d,phantom,sloped,"\simeq" description]\\
\&\widetilde X\times\widetilde X
\end{tikzcd}}
\deleted{est \unsure{transformée}} dans $\mathbb{\Delta}^{\wedge}/\widetilde{X\times X}$, est
$W$-équivalente sur $\widetilde{X\times X}$ à la flèche
\[
\mathrm{Ch}(\widetilde X)\longrightarrow\widetilde X\times\widetilde{\unsure{X|Y}}
\]
de la théorie \unsure{ssimpliciale|simpliciale|semisimpliciale} des chemins (à
supposer qu'il y ait une telle flèche classique,
définie pour \uline{tout} objet $X$ de $\mathbb{\Delta}^{\wedge}$). Ou
qu'on ait \inserted{tout au moins} des isom. canoniques \inserted{dans} $\mathrm{Hot}$
\[
\mathrm{hot}_{\mathrm{Cat}}\bigl(\underline{\mathrm{Ch}}_{\infty}(X;a,b)\bigr)
\simeq\mathrm{hot}_{\mathbb{\Delta}^{\wedge}}\bigl(\mathrm{Ch}(\widetilde X;\tilde a,\tilde b)\bigr).
\]
% archive sheet number: 402
